\documentclass{ximera}

\input{../../../preamble.tex}


\definecolor{fvdnavy}{HTML}{071D43}
\definecolor{fvdcyan}{HTML}{20BFC6}
\definecolor{fvdcyanlight}{HTML}{EFFCFB}
\definecolor{fvdmagenta}{HTML}{F10D69}
\definecolor{fvdmagentalight}{HTML}{FFF1F7}
\definecolor{fvdtext}{HTML}{163044}





%=================================================
% Pedagogische omgevingen
% PDF: tcolorbox
% HTML: learning-box
%=================================================

\newenvironment{voorkennis}
{%
    \ifdefined\HCode
        \HCode{%
<section class="learning-box learning-box--prior">
<div class="learning-box__header">
<span class="learning-box__icon" aria-hidden="true"></span>
<span class="learning-box__title">Voorkennis</span>
</div>
<div class="learning-box__content">
        }%
        \begin{itemize}
    \else
       \begin{tcolorbox}[
    enhanced,
    breakable,
    colback=fvdcyanlight,
    colframe=fvdcyan,
    coltext=fvdtext,
    boxrule=0.5pt,
    leftrule=4pt,
    arc=4mm,
    outer arc=4mm,
    left=7mm,
    right=6mm,
    top=5mm,
    bottom=5mm,
    before skip=6mm,
    after skip=6mm,
    title=Voorkennis,
    coltitle=fvdcyan,
    fonttitle=\bfseries\Large,
    attach title to upper={\par\smallskip},
    boxed title style={
        empty,
        left=0mm,
        right=0mm,
        top=0mm,
        bottom=0mm
    },
    overlay={
        \draw[
            fvdcyan,
            line width=0.5pt
        ]
        ([xshift=42mm,yshift=-13mm]frame.north west)
        --
        ([xshift=-7mm,yshift=-13mm]frame.north east);
    }
]
        \begin{itemize}
    \fi
}
{%
    \end{itemize}
    \ifdefined\HCode
        \HCode{%
</div>
</section>
        }%
    \else
        \end{tcolorbox}
    \fi
}


\newenvironment{leerdoelen}
{%
    \ifdefined\HCode
        \HCode{%
<section class="learning-box learning-box--goals">
<div class="learning-box__header">
<span class="learning-box__icon" aria-hidden="true"></span>
<span class="learning-box__title">Leerdoelen</span>
</div>
<div class="learning-box__content">
        }%
        \begin{itemize}
    \else
\begin{tcolorbox}[
    enhanced,
    breakable,
    colback=fvdmagentalight,
    colframe=fvdmagenta,
    coltext=fvdtext,
    boxrule=0.5pt,
    leftrule=4pt,
    arc=4mm,
    outer arc=4mm,
    left=7mm,
    right=6mm,
    top=5mm,
    bottom=5mm,
    before skip=6mm,
    after skip=6mm,
    title=Leerdoelen,
    coltitle=fvdmagenta,
    fonttitle=\bfseries\Large,
    attach title to upper={\par\smallskip},
    boxed title style={
        empty,
        left=0mm,
        right=0mm,
        top=0mm,
        bottom=0mm
    },
    overlay={
        \draw[
            fvdmagenta,
            line width=0.5pt
        ]
        ([xshift=42mm,yshift=-13mm]frame.north west)
        --
        ([xshift=-7mm,yshift=-13mm]frame.north east);
    }
]
        \begin{itemize}
    \fi
}
{%
    \end{itemize}
    \ifdefined\HCode
        \HCode{%
</div>
</section>
        }%
    \else
        \end{tcolorbox}
    \fi
}


\addPrintStyle{../../..}

\title{Verloop van functies}
\author{Ingmar Herreman}

\begin{abstract}
Een afgeleide is meer dan een nieuwe functie die je kunt berekenen.
Ze vertelt hoe de oorspronkelijke functie verandert.

In dit hoofdstuk brengen we de belangrijkste verbanden uit het vijfde
jaar opnieuw samen. We gebruiken de eerste afgeleide om stijgen,
dalen en extrema te onderzoeken en de tweede afgeleide om de
kromming en buigpunten van een grafiek te bestuderen.

Daarbij staat één idee centraal:
uit informatie over $f'$ en $f''$ leiden we eigenschappen van $f$ af.
\end{abstract}

\begin{document}

% FVD_CHAPTER_DATA_START
% Automatisch gegenereerd uit index.tex.
% Niet handmatig aanpassen.
\fvdchapterdata{1}{Veranderingen in beeld}{3}
% FVD_CHAPTER_DATA_END









































































\maketitle


% ============================================================
% VOORKENNIS EN LEERDOELEN
% ============================================================

\begin{voorkennis}
    \item Je kan veeltermfuncties afleiden.
    \item Je kan een eerste en tweede afgeleide bepalen.
    \item Je kan het teken van een functie onderzoeken.
    \item Je kan nulwaarden algebraïsch bepalen.
    \item Je kan een tekentabel lezen en opstellen.
    \item Je kan het verschil uitleggen tussen $f'(a)$ en $f'(x)$.
  \end{voorkennis}
  
  \begin{leerdoelen}
    \item Je kan het verband leggen tussen het teken van $f'$ en het stijgen of dalen van $f$.
    \item Je kan stijg- en daalintervallen bepalen met behulp van $f'$.
    \item Je kan stationaire punten herkennen en onderzoeken.
    \item Je kan lokale extrema bepalen met behulp van een tekenonderzoek van $f'$.
    \item Je kan uitleggen waarom $f'(a)=0$ op zichzelf niet volstaat om een extremum te besluiten.
    \item Je kan het verband leggen tussen het teken van $f''$ en de kromming van de grafiek.
    \item Je kan buigpunten onderzoeken met behulp van $f''$.
    \item Je kan uitleggen waarom $f''(a)=0$ op zichzelf niet volstaat om een buigpunt te besluiten.
    \item Je kan informatie over $f$, $f'$ en $f''$ met elkaar verbinden.
    \item Je kan eenvoudige wetenschappelijke situaties interpreteren met eerste en tweede afgeleiden.
  \end{leerdoelen}

% ============================================================
% 1 WAAROM LEER JE DIT?
% ============================================================

\FVDsection{Waarom leer je dit?}

In het vorige hoofdstuk hebben we de afgeleide functie opnieuw
opgebouwd. We weten bijvoorbeeld dat

\[
f'(x)
\]

beschrijft hoe $f$ verandert en dat

\[
f''(x)
\]

beschrijft hoe die verandering zelf verandert.

Maar uiteindelijk willen we meestal niet alleen $f'$ en $f''$
\emph{berekenen}. We willen ze gebruiken om iets over $f$ te weten
te komen.

Denk bijvoorbeeld aan de hoogte $h(t)$ van een voorwerp.

Uit

\[
h'(t)>0
\]

kunnen we afleiden dat het voorwerp op dat ogenblik omhoog beweegt.

Uit

\[
h'(t)<0
\]

kunnen we afleiden dat het naar beneden beweegt.

En wanneer de snelheid van positief naar negatief verandert, bereikt
het voorwerp zijn hoogste punt.

We gebruiken dus informatie over een afgeleide om conclusies te
trekken over de oorspronkelijke functie.

\begin{opmerking}
Dat is de belangrijkste denkstap in dit hoofdstuk:

\[
\boxed{
\text{informatie over }f'
\quad\Longrightarrow\quad
\text{informatie over }f
}
\]

en vervolgens

\[
\boxed{
\text{informatie over }f''
\quad\Longrightarrow\quad
\text{informatie over de verandering van }f'
}
\]

Je moet deze verbanden niet alleen kunnen toepassen, maar ook
kunnen verklaren.
\end{opmerking}


% ============================================================
% 2 TEKEN VAN f'
% ============================================================

\FVDsection{Van het teken van $f'$ naar het verloop van $f$}

\FVDsubsection{Wat betekent $f'(x)>0$?}

Neem een interval waarop

\[
f'(x)>0.
\]

In elk punt van dat interval heeft de grafiek van $f$ een raaklijn
met positieve richtingscoëfficiënt.

Als $x$ toeneemt, neemt ook $f(x)$ toe.

De functie is daar dus stijgend.

Op dezelfde manier geldt:

\begin{eigenschap}[Teken van de eerste afgeleide]

Op een interval geldt:

\[
\boxed{
f'(x)>0
\quad\Longrightarrow\quad
f \text{ is stijgend}
}
\]

en

\[
\boxed{
f'(x)<0
\quad\Longrightarrow\quad
f \text{ is dalend}.
}
\]

\end{eigenschap}

Dit is een uitspraak over een \emph{interval}.

Eén enkele waarde zoals

\[
f'(2)>0
\]

vertelt alleen wat er in de buurt van $x=2$ gebeurt. Daaruit mag je
niet besluiten dat $f$ op haar volledige domein stijgend is.


% ============================================================
% EERSTE KORTE CHECK
% ============================================================

\begin{oefening}[Van afgeleide naar verloop]{\oefB\ESS}

Van een functie $f$ weten we:

\[
f'(x)>0
\qquad\text{voor }-3<x<2
\]

en

\[
f'(x)<0
\qquad\text{voor }x>2.
\]

Vul aan.

\begin{question}
Op $]-3,2[$ is $f$

\begin{multipleChoice}
\choice[correct]{stijgend.}
\choice{dalend.}
\choice{constant.}
\end{multipleChoice}
\end{question}

\begin{question}
Op $]2,+\infty[$ is $f$

\begin{multipleChoice}
\choice{stijgend.}
\choice[correct]{dalend.}
\choice{constant.}
\end{multipleChoice}
\end{question}

\end{oefening}


% ============================================================
% 3 TEKENONDERZOEK
% ============================================================

\FVDsection{Stijgen en dalen bepalen}

Om te onderzoeken waar een functie stijgt of daalt, kunnen we dus
het teken van $f'$ onderzoeken.

De werkwijze is:

\[
f
\longrightarrow
f'
\longrightarrow
\text{nulwaarden van }f'
\longrightarrow
\text{teken van }f'
\longrightarrow
\text{verloop van }f.
\]

\begin{voorbeeld}[Stijgen en dalen van een derdegraadsfunctie]

Gegeven:

\[
f(x)=x^3-3x^2-9x+5.
\]

We onderzoeken waar $f$ stijgt en daalt.

\textbf{Stap 1: bepaal de afgeleide.}

\[
f'(x)=3x^2-6x-9.
\]

Factoriseren geeft

\[
f'(x)=3(x^2-2x-3)
\]

en dus

\[
\boxed{
f'(x)=3(x+1)(x-3).
}
\]

\textbf{Stap 2: bepaal de nulwaarden van $f'$.}

\[
f'(x)=0
\]

geeft

\[
x=-1
\qquad\text{of}\qquad
x=3.
\]

Deze waarden verdelen de getallenas in drie intervallen.

\[
]-\infty,-1[,
\qquad
]-1,3[,
\qquad
]3,+\infty[.
\]

\textbf{Stap 3: onderzoek het teken van $f'$.}

Omdat

\[
f'(x)=3(x+1)(x-3),
\]

vinden we:

\[
\begin{array}{c|ccccc}
x
& ]-\infty,-1[
& -1
& ]-1,3[
& 3
& ]3,+\infty[
\\ \hline
f'(x)
& +
& 0
& -
& 0
& +
\end{array}
\]

\textbf{Stap 4: vertaal naar het verloop van $f$.}

Dus:

\[
\boxed{
f \text{ stijgt op }]-\infty,-1[
\text{ en } ]3,+\infty[
}
\]

en

\[
\boxed{
f \text{ daalt op }]-1,3[.
}
\]

\end{voorbeeld}


% ============================================================
% 4 STATIONAIRE PUNTEN
% ============================================================

\FVDsection{Stationaire punten}

De waarden waarvoor

\[
f'(x)=0
\]

verdienen bijzondere aandacht.

\begin{definitie}[Stationair punt]

Een punt

\[
P(a,f(a))
\]

van de grafiek van een afleidbare functie $f$ heet een
\textbf{stationair punt} als

\[
\boxed{f'(a)=0.}
\]

De raaklijn aan de grafiek is daar horizontaal.

\end{definitie}

Maar een stationair punt is niet noodzakelijk een maximum of minimum.

Dat hebben we in hoofdstuk 1 al aangekondigd.

Vergelijk bijvoorbeeld:

\[
f(x)=x^2
\]

en

\[
g(x)=x^3.
\]

Voor beide functies geldt:

\[
f'(0)=0
\qquad\text{en}\qquad
g'(0)=0.
\]

Toch heeft $f$ in $x=0$ een minimum, terwijl $g$ daar geen extremum
heeft.

We moeten dus onderzoeken wat het teken van de afgeleide
\emph{voor en na} het stationaire punt doet.


% ============================================================
% 5 EXTREMA
% ============================================================

\FVDsection{Extrema herkennen met $f'$}

\FVDsubsection{Van stijgen naar dalen}

Stel dat bij $x=a$ het teken van $f'$ verandert als

\[
+\quad 0\quad -.
\]

Links van $a$ stijgt $f$.

Rechts van $a$ daalt $f$.

De functie gaat dus eerst omhoog en daarna omlaag.

Daarom heeft $f$ in $x=a$ een lokaal maximum.

\FVDsubsection{Van dalen naar stijgen}

Wanneer het teken verandert als

\[
-\quad 0\quad +,
\]

daalt $f$ eerst en stijgt ze daarna.

Dan heeft $f$ in $x=a$ een lokaal minimum.

\begin{eigenschap}[Extrema via het teken van $f'$]

Als $f'$ bij $x=a$ van teken verandert, dan geldt:

\[
\boxed{
+\;\longrightarrow\;-
\quad\Rightarrow\quad
\text{lokaal maximum}
}
\]

en

\[
\boxed{
-\;\longrightarrow\;+
\quad\Rightarrow\quad
\text{lokaal minimum}.
}
\]

Als $f'(a)=0$ maar $f'$ niet van teken verandert, kunnen we op basis
daarvan geen lokaal extremum besluiten.

\end{eigenschap}


% ============================================================
% VOORBEELD EXTREMA
% ============================================================

\begin{voorbeeld}[De extrema volledig bepalen]

We keren terug naar

\[
f(x)=x^3-3x^2-9x+5.
\]

We vonden:

\[
\begin{array}{c|ccccc}
x
& ]-\infty,-1[
& -1
& ]-1,3[
& 3
& ]3,+\infty[
\\ \hline
f'(x)
& +
& 0
& -
& 0
& +
\\ \hline
f
& \nearrow
&
& \searrow
&
& \nearrow
\end{array}
\]

Bij $x=-1$ verandert $f'$ van positief naar negatief.

Dus heeft $f$ daar een lokaal maximum.

De functiewaarde is

\[
f(-1)
=
(-1)^3-3(-1)^2-9(-1)+5
=
10.
\]

Het lokale maximumpunt is dus

\[
\boxed{(-1,10)}.
\]

Bij $x=3$ verandert $f'$ van negatief naar positief.

Dus heeft $f$ daar een lokaal minimum.

\[
f(3)
=
27-27-27+5
=
-22.
\]

Het lokale minimumpunt is

\[
\boxed{(3,-22)}.
\]

\end{voorbeeld}


% ============================================================
% OEFENING EXTREMA
% ============================================================

\begin{oefening}[Lees de tekentabel]{\oefB\ESS}

Van een afleidbare functie $f$ is gegeven:

\[
\begin{array}{c|ccccc}
x
& ]-\infty,1[
& 1
& ]1,4[
& 4
& ]4,+\infty[
\\ \hline
f'(x)
& -
& 0
& +
& 0
& -
\end{array}
\]

\begin{question}
Welke uitspraak over $x=1$ is correct?

\begin{multipleChoice}
\choice{Daar ligt noodzakelijk een lokaal maximum.}
\choice[correct]{Daar ligt een lokaal minimum.}
\choice{Daar ligt geen stationair punt.}
\end{multipleChoice}
\end{question}

\begin{question}
Welke uitspraak over $x=4$ is correct?

\begin{multipleChoice}
\choice[correct]{Daar ligt een lokaal maximum.}
\choice{Daar ligt een lokaal minimum.}
\choice{De functie is daar noodzakelijk niet afleidbaar.}
\end{multipleChoice}
\end{question}

\end{oefening}


% ============================================================
% 6 BELANGRIJKE DENKFOUT
% ============================================================

\FVDsection{Een nulwaarde van $f'$ is niet genoeg}

Een veelgemaakte redenering is:

\[
f'(a)=0
\quad\Rightarrow\quad
f \text{ heeft een extremum in }a.
\]

Die redenering is fout.

Neem

\[
f(x)=x^3.
\]

Dan

\[
f'(x)=3x^2
\]

en dus

\[
f'(0)=0.
\]

Maar

\[
f'(x)>0
\]

zowel voor $x<0$ als voor $x>0$.

Het tekenverloop is:

\[
+\quad 0\quad +.
\]

De functie blijft dus stijgen.

Er is geen maximum en geen minimum.

\begin{opmerking}
De vergelijking

\[
f'(x)=0
\]

levert \emph{kandidaten} voor stationaire extrema.

Daarna moet je onderzoeken wat er met het teken van $f'$ gebeurt.
\end{opmerking}


\begin{oefening}[Is er een extremum?]{\oefU\ESS}

Een leerling onderzoekt

\[
f(x)=x^3
\]

en schrijft:

\begin{quote}
Omdat $f'(0)=0$, heeft $f$ in $x=0$ een extremum.
\end{quote}

Leg uit waarom deze redenering niet geldig is en geef het correcte
besluit.

\end{oefening}


% ============================================================
% 7 TWEEDE AFGELEIDE
% ============================================================

\FVDsection{Wat vertelt de tweede afgeleide?}

Tot nu toe gebruikten we $f'$ om te onderzoeken of $f$ stijgt of
daalt.

Maar ook $f'$ zelf kan stijgen of dalen.

Daarover geeft $f''$ informatie.

Herinner je:

\[
f''=(f')'.
\]

Dus:

\[
f''(x)>0
\]

betekent dat $f'$ stijgt, terwijl

\[
f''(x)<0
\]

betekent dat $f'$ daalt.

Dit zegt iets over de manier waarop de hellingen van de grafiek van
$f$ veranderen.

\begin{eigenschap}[Teken van de tweede afgeleide]

Op een interval geldt:

\[
\boxed{
f''(x)>0
\quad\Longrightarrow\quad
f' \text{ is stijgend}
}
\]

en

\[
\boxed{
f''(x)<0
\quad\Longrightarrow\quad
f' \text{ is dalend}.
}
\]

Daardoor geeft $f''$ informatie over de \textbf{kromming} van de
grafiek van $f$.

\end{eigenschap}


% ============================================================
% 8 KROMMING
% ============================================================

\FVDsection{Kromming van de grafiek}

Bekijk eerst een stijgende functie.

Het is mogelijk dat de functie stijgt en tegelijk steeds steiler
wordt.

Dan neemt $f'$ toe en geldt typisch

\[
f''>0.
\]

Maar een stijgende functie kan ook steeds minder steil worden.

Dan neemt $f'$ af en geldt typisch

\[
f''<0.
\]

Daarom zijn ``stijgend'' en ``positieve tweede afgeleide'' twee
verschillende eigenschappen.

\begin{opmerking}
Verwar de rollen van de twee afgeleiden niet:

\[
\boxed{
f'
\text{ vertelt vooral of }f\text{ stijgt of daalt}
}
\]

terwijl

\[
\boxed{
f''
\text{ vertelt hoe de helling van }f\text{ verandert}.
}
\]

Een functie kan dus tegelijk stijgend zijn en $f''<0$ hebben.
\end{opmerking}


% ============================================================
% HOL/BOL
% ============================================================

\FVDsubsection{Hol en bol}

In deze cursus gebruiken we de volgende terminologie:

\begin{eigenschap}[Kromming]

Op een interval:

\[
f''(x)>0
\]

betekent dat de hellingen toenemen.

De grafiek is daar \textbf{hol}.

Wanneer

\[
f''(x)<0,
\]

nemen de hellingen af.

De grafiek is daar \textbf{bol}.

\end{eigenschap}

\begin{opmerking}
De woorden \emph{hol} en \emph{bol} worden in verschillende bronnen
niet altijd op dezelfde manier gebruikt.

Daarom is het veiliger om bij een redenering ook het teken van
$f''$ te vermelden.
\end{opmerking}


% ============================================================
% OEFENING f''
% ============================================================

\begin{oefening}[Wat kun je besluiten uit $f''$?]{\oefB\ESS}

Op een interval $I$ geldt

\[
f''(x)>0.
\]

Welke uitspraak volgt daaruit?

\begin{multipleChoice}
\choice{$f$ is noodzakelijk stijgend op $I$.}
\choice{$f$ is noodzakelijk positief op $I$.}
\choice[correct]{$f'$ is stijgend op $I$.}
\choice{$f'$ is negatief op $I$.}
\end{multipleChoice}

\end{oefening}


% ============================================================
% 9 BUIGPUNT
% ============================================================

\FVDsection{Buigpunten}

Een buigpunt is een punt waar de aard van de kromming verandert.

\begin{definitie}[Buigpunt]

Een punt van de grafiek van $f$ waar de kromming verandert, noemen
we een \textbf{buigpunt}.

Bij een tweemaal afleidbare functie onderzoeken we daarvoor het
teken van $f''$.

\end{definitie}

Wanneer $f''$ bijvoorbeeld verandert als

\[
-\quad 0\quad +,
\]

verandert de kromming.

Het overeenkomstige punt van de grafiek is dan een buigpunt.


% ============================================================
% DENKFOUT f''
% ============================================================

\FVDsection{Ook $f''(a)=0$ is niet genoeg}

Er ontstaat nu een situatie die sterk lijkt op die bij extrema.

Voor extrema was

\[
f'(a)=0
\]

op zichzelf niet voldoende.

Voor buigpunten is

\[
f''(a)=0
\]

op zichzelf evenmin voldoende.

Neem bijvoorbeeld

\[
f(x)=x^4.
\]

Dan

\[
f'(x)=4x^3
\]

en

\[
f''(x)=12x^2.
\]

We hebben inderdaad

\[
f''(0)=0.
\]

Maar voor $x\neq0$ geldt

\[
12x^2>0.
\]

Het teken van $f''$ is dus:

\[
+\quad 0\quad +.
\]

Er is geen verandering van kromming.

Dus is $(0,0)$ geen buigpunt.

\begin{opmerking}
Ook hier geldt:

\[
f''(a)=0
\]

geeft een mogelijke kandidaat, maar je moet onderzoeken of de
kromming werkelijk verandert.
\end{opmerking}


% ============================================================
% VOORBEELD BUIGPUNT
% ============================================================

\begin{voorbeeld}[Een buigpunt bepalen]

Neem opnieuw

\[
f(x)=x^3-3x^2-9x+5.
\]

We vonden:

\[
f'(x)=3x^2-6x-9.
\]

Dus:

\[
f''(x)=6x-6.
\]

Een mogelijke verandering van kromming vinden we uit

\[
f''(x)=0.
\]

Dus:

\[
6x-6=0
\]

en bijgevolg

\[
x=1.
\]

Voor $x<1$ geldt

\[
f''(x)<0,
\]

en voor $x>1$:

\[
f''(x)>0.
\]

Het teken verandert dus:

\[
-\quad 0\quad +.
\]

Er is werkelijk een verandering van kromming.

De functiewaarde is

\[
f(1)=1-3-9+5=-6.
\]

Het buigpunt is dus

\[
\boxed{(1,-6)}.
\]

\end{voorbeeld}


% ============================================================
% OEFENING BUIGPUNT
% ============================================================

\begin{oefening}[Buigpunt of niet?]{\oefU\ESS}

Onderzoek voor elke functie of de oorsprong een buigpunt is.

\begin{question}

\[
f(x)=x^3
\]

Bepaal $f''(x)$ en onderzoek het teken links en rechts van $0$.

\end{question}

\begin{question}

\[
g(x)=x^4
\]

Bepaal $g''(x)$ en onderzoek het teken links en rechts van $0$.

\end{question}

Leg daarna uit waarom

\[
f''(0)=0
\]

niet voldoende is om te besluiten dat $(0,f(0))$ een buigpunt is.

\end{oefening}


% ============================================================
% 10 ALLES SAMEN
% ============================================================

\FVDsection{$f$, $f'$ en $f''$ samen lezen}

We kunnen de belangrijkste verbanden nu samenbrengen.

\[
\boxed{
f'(x)>0
\quad\Longrightarrow\quad
f \text{ stijgt}
}
\]

\[
\boxed{
f'(x)<0
\quad\Longrightarrow\quad
f \text{ daalt}
}
\]

\[
\boxed{
f' : +\to-
\quad\Longrightarrow\quad
\text{lokaal maximum}
}
\]

\[
\boxed{
f' : -\to+
\quad\Longrightarrow\quad
\text{lokaal minimum}
}
\]

en

\[
\boxed{
f''(x)>0
\quad\Longrightarrow\quad
f' \text{ stijgt}
}
\]

\[
\boxed{
f''(x)<0
\quad\Longrightarrow\quad
f' \text{ daalt}.
}
\]

Een tekenwisseling van $f''$ wijst op een verandering van kromming
en dus op een buigpunt.


% ============================================================
% CONCEPTUELE OEFENING
% ============================================================

\begin{oefening}[Van $f'$ naar $f$]{\oefU\ESS}

Van een functie $f$ is de volgende informatie over $f'$ bekend:

\[
\begin{array}{c|ccccc}
x
& ]-\infty,-2[
& -2
& ]-2,3[
& 3
& ]3,+\infty[
\\ \hline
f'(x)
& +
& 0
& -
& 0
& +
\end{array}
\]

Zonder een voorschrift van $f$ te kennen:

\begin{question}
Bepaal de intervallen waarop $f$ stijgt.
\end{question}

\begin{question}
Bepaal de intervallen waarop $f$ daalt.
\end{question}

\begin{question}
Classificeer de stationaire punten bij $x=-2$ en $x=3$.
\end{question}

\begin{question}
Kun je uit deze gegevens bepalen waar de buigpunten liggen?
Verantwoord.
\end{question}

\end{oefening}


% ============================================================
% 11 FUNCTIEONDERZOEK
% ============================================================

\FVDsection{Een functie systematisch onderzoeken}

Voor een veeltermfunctie kunnen we de informatie nu systematisch
organiseren.

Een mogelijke werkwijze is:

\begin{enumerate}

\item Bepaal $f'(x)$.

\item Los
\[
f'(x)=0
\]
op.

\item Onderzoek het teken van $f'$.

\item Besluit waar $f$ stijgt en daalt.

\item Onderzoek eventuele extrema en bereken hun coördinaten.

\item Bepaal $f''(x)$.

\item Onderzoek het teken van $f''$.

\item Bepaal waar de kromming verandert.

\item Bereken de coördinaten van eventuele buigpunten.

\item Combineer alle informatie tot één coherent beeld van de functie.

\end{enumerate}

\begin{opmerking}
Een functieonderzoek is geen verzameling losse berekeningen.

Elke stap moet informatie opleveren over de grafiek.

Vraag jezelf daarom telkens af:

\[
\boxed{\text{Wat vertelt dit resultaat mij over }f?}
\]
\end{opmerking}


% ============================================================
% VOLLEDIGE OEFENING
% ============================================================

\begin{oefening}[Volledig verlooponderzoek]{\oefU\ESS}

Gegeven:

\[
f(x)=x^3-6x^2+9x+1.
\]

Onderzoek het verloop van $f$.

\begin{question}
Bepaal $f'(x)$ en factoriseer zo ver mogelijk.
\end{question}

\begin{question}
Bepaal de stationaire punten.
\end{question}

\begin{question}
Stel een tekentabel op voor $f'$ en bepaal waar $f$ stijgt en daalt.
\end{question}

\begin{question}
Bepaal en classificeer de lokale extrema.
\end{question}

\begin{question}
Bepaal $f''(x)$.
\end{question}

\begin{question}
Onderzoek de kromming van de grafiek en bepaal het buigpunt.
\end{question}

\begin{question}
Vat je resultaten samen in één verloopschema.
\end{question}

\end{oefening}


% ============================================================
% 12 WETENSCHAPPELIJKE CONTEXT
% ============================================================

\FVDsection{Van positie naar snelheid en versnelling}

Een bijzonder belangrijke toepassing vinden we in de mechanica.

Stel dat

\[
s(t)
\]

de positie van een bewegend voorwerp beschrijft.

Dan is

\[
v(t)=s'(t)
\]

de ogenblikkelijke snelheid.

De afgeleide daarvan is

\[
a(t)=v'(t)=s''(t),
\]

de versnelling.

We krijgen dus:

\[
\boxed{
s
\longrightarrow
s'=v
\longrightarrow
s''=v'=a
}
\]

Dit is precies dezelfde structuur als

\[
f
\longrightarrow
f'
\longrightarrow
f''.
\]


% ============================================================
% STEM-OEFENING
% ============================================================

\begin{oefening}[Beweging interpreteren]{\oefU\ESS}

De positie van een bewegend voorwerp langs een rechte wordt gegeven
door

\[
s(t)=t^3-6t^2+9t,
\qquad t\geq0,
\]

waarbij $s$ in meter en $t$ in seconde wordt uitgedrukt.

\begin{question}
Bepaal de snelheid $v(t)$.
\end{question}

\begin{question}
Op welke tijdstippen staat het voorwerp ogenblikkelijk stil?
\end{question}

\begin{question}
Onderzoek op welke tijdsintervallen het voorwerp in positieve en in
negatieve richting beweegt.
\end{question}

\begin{question}
Bepaal de versnelling $a(t)$.
\end{question}

\begin{question}
Op welk tijdstip verandert de versnelling van teken?
Wat betekent dat voor de snelheid?
\end{question}

\end{oefening}


% ============================================================
% 13 REDENEREN ZONDER REKENEN
% ============================================================

\FVDsection{Redeneren zonder voorschrift}

Bij analyseren krijg je niet altijd een functievoorschrift.

Soms krijg je alleen informatie over een functie of haar afgeleiden.

\begin{oefening}[Wat moet waar zijn?]{\oefU\ESS}

Een tweemaal afleidbare functie $f$ voldoet op een interval $I$ aan

\[
f'(x)>0
\qquad\text{en}\qquad
f''(x)<0.
\]

Welke beschrijving is correct?

\begin{multipleChoice}
\choice{$f$ daalt en wordt steeds steiler.}
\choice{$f$ daalt en wordt steeds vlakker.}
\choice{$f$ stijgt en wordt steeds steiler.}
\choice[correct]{$f$ stijgt, maar haar helling neemt af.}
\end{multipleChoice}

Leg je keuze uit met behulp van de betekenis van $f'$ en $f''$.

\end{oefening}


% ============================================================
% 14 OMGEKEERD REDENEREN
% ============================================================

\FVDsection{Van het gedrag van $f$ naar haar afgeleiden}

Tot nu toe redeneerden we meestal van de afgeleiden naar de functie.

Maar je moet ook in de andere richting kunnen denken.

Als $f$ stijgt, verwachten we:

\[
f'>0.
\]

Als de grafiek tijdens dat stijgen steeds steiler wordt, neemt $f'$
toe en verwachten we:

\[
f''>0.
\]

Als de grafiek stijgt maar steeds vlakker wordt, blijft

\[
f'>0,
\]

terwijl

\[
f''<0.
\]

Dat onderscheid is belangrijk.

\begin{oefening}[Van functie naar afgeleiden]{\oefU\ESS}

Een functie $f$ is op een interval stijgend. Haar grafiek wordt van
links naar rechts steeds vlakker.

Welke combinatie past daarbij?

\begin{multipleChoice}
\choice{$f'<0$ en $f''<0$}
\choice{$f'<0$ en $f''>0$}
\choice{$f'>0$ en $f''>0$}
\choice[correct]{$f'>0$ en $f''<0$}
\end{multipleChoice}

\end{oefening}


% ============================================================
% 15 VEELGEMAAKTE DENKFOUTEN
% ============================================================

\FVDsection{Veelgemaakte denkfouten}

\begin{opmerking}

\textbf{Denkfout 1}

\[
f'(a)>0
\]

betekent niet dat $f$ overal stijgt.

Het zegt iets over de verandering rond $x=a$.

\medskip

\textbf{Denkfout 2}

\[
f'(a)=0
\]

betekent niet automatisch dat $f$ een extremum heeft.

Onderzoek het teken van $f'$.

\medskip

\textbf{Denkfout 3}

\[
f''(a)=0
\]

betekent niet automatisch dat $f$ een buigpunt heeft.

Onderzoek of de kromming verandert.

\medskip

\textbf{Denkfout 4}

\[
f''>0
\]

betekent niet dat $f$ stijgt.

Het betekent dat $f'$ stijgt.

\medskip

\textbf{Denkfout 5}

Een extremum is een \emph{punt} van de grafiek.

Wanneer $x=a$ de extremumplaats is, is

\[
(a,f(a))
\]

het bijbehorende extremumpunt.

\end{opmerking}


% ============================================================
% 16 VERBIND ALLES
% ============================================================

\FVDsection{Kun je alles verbinden?}

\begin{oefening}[Eén functie, drie niveaus]{\oefG}

Een functie $f$ is tweemaal afleidbaar.

Er geldt:

\[
f'(2)=0,
\qquad
f''(2)=0.
\]

Een leerling beweert:

\begin{quote}
Bij $x=2$ ligt dus zowel een extremum als een buigpunt.
\end{quote}

Onderzoek deze uitspraak kritisch.

Leg uit:
\begin{itemize}
  \item waarom de gegeven informatie niet voldoende is;
  \item welke bijkomende informatie over $f'$ nodig is om een extremum
        te onderzoeken;
  \item welke bijkomende informatie over $f''$ nodig is om een buigpunt
        te onderzoeken;
  \item welke verschillende situaties mogelijk zijn.
\end{itemize}

\end{oefening}


% ============================================================
% 17 SAMENVATTING
% ============================================================

\FVDsection{Samengevat}

De eerste afgeleide beschrijft de verandering van $f$:

\[
f'(x)>0
\Rightarrow
f \text{ stijgt},
\]

\[
f'(x)<0
\Rightarrow
f \text{ daalt}.
\]

Een tekenwisseling van $f'$ kan een extremum opleveren:

\[
+\to-
\Rightarrow
\text{lokaal maximum},
\]

\[
-\to+
\Rightarrow
\text{lokaal minimum}.
\]

De tweede afgeleide beschrijft de verandering van $f'$:

\[
f''(x)>0
\Rightarrow
f' \text{ stijgt},
\]

\[
f''(x)<0
\Rightarrow
f' \text{ daalt}.
\]

Een verandering van teken van $f''$ wijst op een verandering van
kromming.

Het belangrijkste verband van dit hoofdstuk is daarom:

\[
\boxed{
f
\;\longrightarrow\;
f'
\;\longrightarrow\;
f''
}
\]

maar bij een analyse redeneren we ook terug:

\[
\boxed{
f''
\;\longrightarrow\;
\text{gedrag van }f'
\;\longrightarrow\;
\text{gedrag van }f.
}
\]


% ============================================================
% 18 MINIMUMTRAJECT
% ============================================================

\FVDsection{Wat moet je zeker beheersen?}

Na deze opfrissing moet je niet alleen eerste en tweede afgeleiden
kunnen berekenen.

Je moet informatie uit $f'$ en $f''$ kunnen \textbf{analyseren} en
vertalen naar eigenschappen van $f$.

Je moet zelfstandig kunnen:

\begin{itemize}

\item uit het teken van $f'$ besluiten waar $f$ stijgt of daalt;

\item met een tekenonderzoek van $f'$ lokale extrema herkennen;

\item uitleggen waarom $f'(a)=0$ niet automatisch een extremum geeft;

\item uit het teken van $f''$ afleiden hoe de hellingen van $f$
veranderen;

\item een verandering van kromming onderzoeken;

\item uitleggen waarom $f''(a)=0$ niet automatisch een buigpunt geeft;

\item informatie over $f$, $f'$ en $f''$ in beide richtingen met elkaar
verbinden;

\item deze begrippen betekenisvol interpreteren in een eenvoudige
wetenschappelijke context.

\end{itemize}

\begin{opmerking}
De oefeningen met de diamant vormen in dit hoofdstuk het
minimumtraject.

Concreet zijn dat:

\begin{itemize}
  \item \textit{Van afgeleide naar verloop};
  \item \textit{Lees de tekentabel};
  \item \textit{Is er een extremum?};
  \item \textit{Wat kun je besluiten uit $f''$?};
  \item \textit{Buigpunt of niet?};
  \item \textit{Van $f'$ naar $f$};
  \item \textit{Volledig verlooponderzoek};
  \item \textit{Beweging interpreteren};
  \item \textit{Wat moet waar zijn?};
  \item \textit{Van functie naar afgeleiden}.
\end{itemize}

Merk op dat verschillende essentiële oefeningen het label
\oefU{} dragen. Dat is bewust: het leerplandoel vraagt uiteindelijk
\emph{analyseren}. Alleen standaardtechnieken uitvoeren is daarvoor
niet voldoende.
\end{opmerking}


% ============================================================
% OVERGANG
% ============================================================

\FVDsection{Van opfrissing naar nieuwe analyse}

We hebben nu de belangrijkste ideeën uit het vijfde jaar opnieuw
samengebracht:

\[
\text{verandering}
\longrightarrow
\text{afgeleid getal}
\longrightarrow
\text{afgeleide functie}
\longrightarrow
\text{verloop van een functie}.
\]

Tot nu toe werkten we vooral met veeltermfuncties waarvoor de
afgeleiden relatief eenvoudig te bepalen zijn.

In de volgende thema's breiden we onze gereedschapskist uit.

We zullen dezelfde ideeën gebruiken bij andere functieklassen en
nieuwe afleidingstechnieken. Daardoor wordt het mogelijk om veel
rijkere functies te analyseren en later ook extremumproblemen
systematisch aan te pakken.

\end{document}